\subsection{Actual higher-weight occurrence}
\label{sec:selmer-arithmetic-point}

The arithmetic-point implication can be proved directly on the original
ordinary modules.  The only higher-weight comparison needed is
nonvanishing of lowest-weight evaluation on a classical ordinary
vector.

Fix the original Hida datum: an odd prime \(\ell\), a CM extension
\(L/F\) unramified at every finite place, its definite unitary group,
and the original coefficient ring, residual characters, tame type,
weight ring \(A\), universal ordinary determinant ring \(B\), and
residual ordinary Hecke module \(M_A\).  Assume \(S'=S\).  The exact
equality \(S'=S\cup\operatorname{image}(Q)\), with \(Q\) outside
\(S\), then makes the auxiliary Taylor--Wiles set empty.  Keep the
original principal level at \(v_0\).

\begin{proposition}[The full original arithmetic-point implication]
\label{prop:selmer-arithmetic-point}
Let \(J\) be the original Hecke ideal: either \(M_A=0,J=B\), or
\(J=\ker f\) for the prescribed surjection
\(f:B\twoheadrightarrow T\subseteq\operatorname{End}_A(M_A)\)
identifying the universal determinant and ordered characters with the
actual Hecke determinant and characters.

For every \(\mathcal O\)-algebra point
\(x:B\to\overline{\mathbb Q}_\ell\), let \(\rho\) be a continuous
semisimple representation of rank \(2\leq n\leq4\) whose determinant
is \(xD_B\).  Suppose \(\rho\) is attached in the original normalization
to a regular algebraic conjugate self-dual cuspidal representation,
is unramified at the prescribed place above \(v_0\), and has at every
chosen \(w\mid\ell\) the ordered flag with characters
\(\psi_i=x\psi_{B,i}\).  Assume these characters are locally algebraic
of integral labeled weights
\[
 h_{0,\tau}<\cdots<h_{n-1,\tau},\qquad
                  \psi_i\ne\psi_j\epsilon\quad\hbox{for all }i,j.
                                                        \tag{AP1}
\]
Then \(M_A\ne0\) and \(J\subseteq\ker x\).  This supplies the entire
original \(\mathrm{PT}\) contract, without an additional distinctness
hypothesis on finite inertial characters.
\end{proposition}

\begin{proof}
\emph{Points and descent.}
Algebraic points of \(B\) are automatically integral and continuous.
Indeed, if \(b\in\mathfrak m_B\) and \(|x(b)|_\ell\geq1\), a monic
integral equation for \(x(b)^{-1}\) makes
\(1+c_1b+\cdots+c_db^d\), a unit of \(B\), map to zero.
Thus \(x(\mathfrak m_B)\) has positive valuation.  Every element of
\(B\) is a residue lift from \(\mathcal O\) plus an element of
\(\mathfrak m_B\), so all coefficient images are integral.  Finite generators
of \(\mathfrak m_B\) give a positive lower valuation bound on their
nonzero images, proving continuity.  Their images generate a finite
extension \(E\) containing \(x(B)\), by completeness and polynomial
approximation.  Coefficients may henceforth be enlarged inside
\(\overline{\mathbb Q}_\ell\).

The universal trace relations and nondegeneracy of the trace pairing
on the image algebra of a semisimple characteristic-zero representation
make \(\rho\) factor through the original \(G_0\).  Thus it is
unramified outside \(S\).  The adopted classical local/global theorem
makes its cuspidal representation \(\Pi\) spherical there and at
\(v_0\).  The original unitary group is quasi-split at every finite
place: at an inert place choose a unit of norm \(-1\) in the unramified
quadratic extension and split the standard form into hyperbolic planes
and at most one line.  The degree \([F:\mathbb Q]\) is even: the quadratic idele character
of \(L/F\) is one on \(-1\) at finite places and minus one at real
places, so its value on the principal idele \(-1\) gives the parity.
Thus the definite even-rank parity and the required degree at least
two both hold over the original field.

Labesse's Theorem~5.4 gives a descent \(\sigma\) to this definite group
\cite{SelmerPTLabesse}.  Its finite-place unramifiedness and quasi-split
hypotheses have just been checked.  In the proof, spherical test
functions at inert places retain the nonzero twisted trace of \(\Pi\)
by the unramified fundamental lemma; thus the descent may be chosen
spherical at all those places.  Compact infinity gives discrete series
and property \((*)\); Theorem~5.9 gives strong split-place base change
and the matching infinitesimal character.  No multiplicity-one
assertion is needed.

At a prescribed type place the universal relation \(\mathrm{Pair.I}\)
and the same trace pairing kill
\(I_u\cap\bigcap_i\ker\bar\chi_i\).  The remaining inertia quotient
is finite abelian of order prime to \(\ell\).  Its residual determinant
therefore identifies the exact Teichm\"uller characters and their
multiplicities, with zero monodromy.  The established exact type theorem
supplies a vector for the original tame lattice.  At \(v_0\) use the
spherical vector, which is fixed by the original principal level.

\emph{The local principal series.}
For \(k=L_w\), let \(q\) be its residue cardinality and use the original
Artin map: \(v_k(\operatorname{art}(g))=-\deg(g)\), where
arithmetic Frobenius has degree one.  Thus \(g_\varpi\) has degree
minus one.
The locally algebraic characters give smooth \(\omega_i:k^\times\to
\mathbb C^\times\) satisfying
\[
 \psi_i(g_a)=\iota^{-1}\omega_i(a)
       \prod_{\tau:k\hookrightarrow E}\tau(a)^{-h_{i,\tau}},
 \qquad
 \epsilon(g_a)=|a|_kN_{k/\mathbb Q_\ell}(a).             \tag{AP2}
\]
The second identity is the source's exact \texttt{epsArt} formula
\(\epsilon(g)=q^{\deg(g)}N(\operatorname{art}(g))\); its
\texttt{UnifVal} calculation explicitly gives
\(\epsilon(g_\varpi)=N(\varpi)/q\)~\cite{SelmerPTSource}.
The rank-one de Rham dictionary, exactness on the given flag, and
semisimplified local compatibility \cite[Theorem~1.1]{SelmerPTBLGGTLocal}
give
\(\operatorname{rec}(\Pi_w|\det|^{(1-n)/2})^{\mathrm{ss}}
=\bigoplus_i\omega_i\).
Every \(\psi_i(g_a)\) is an \(\ell\)-adic unit by compactness.
Equality \(\omega_i=\omega_j\) would therefore force
\(\sum_\tau(h_{i,\tau}-h_{j,\tau})=0\), contrary to strict ordering.
If \(\omega_i/\omega_j=|\cdot|_k\), valuation at \(\varpi\) gives
\[
 \sum_\tau(h_{j,\tau}-h_{i,\tau})=[k:\mathbb Q_\ell].    \tag{AP3}
\]
All summands have the same sign and are integers.  Hence \(j=i+1\)
and each difference is one.  Equation~\textup{(AP2)} then gives
\(\psi_i=\psi_j\epsilon\), contradicting \textup{(AP1)}.
Interchanging the indices excludes the inverse norm ratio.
Monodromy would join characters with one of these ratios, so it is
zero.  Principal-series irreducibility \cite[Theorem~4.2]{SelmerPTBZ}
and the exact local correspondence give
\[
 \sigma_w\simeq\Pi_w\simeq
 \operatorname{Ind}_{B(k)}^{\mathrm{GL}_n(k),\mathrm{norm}}
       (\omega_i|\cdot|_k^{(n-1)/2})_{0\leq i<n}.        \tag{AP4}
\]
This argument requires neither Shin-regularity nor a preliminary
Iwahori condition.

\emph{The integral coefficient representation.}
Set \(\lambda_{\tau,n-1-i}=h_{i,\tau}-i\); these are dominant
integral weights.  The exact Hodge--Tate formula
\cite[Theorem~2.1.1]{SelmerPTBLGGT} identifies the infinity coefficient
of \(\sigma\) with \(\lambda^\vee\), including the reversal and sign.
Construct its dual coefficient lattice as follows.  For one embedding,
put \(d_i=\lambda_{i-1}-\lambda_i\) and take
\[
 \mathcal T_\lambda=
 \bigotimes_{i=1}^{n-1}(\bigwedge^i\mathcal O_E^n)^{\otimes d_i}
                          \otimes\det^{\lambda_{n-1}}.
                                                        \tag{AP5}
\]
Intersect with the characteristic-zero Cartan component of highest
weight \(\lambda\).  The resulting lattice is saturated, finite free,
and stable under the compact integral group.  Its unique lowest-weight
tensor is primitive; decompose by the ambient integral weight spaces
and choose it as the first basis vector.  Every other weight differs
from it by a nonzero sum of positive roots.  Tensor over the original
chosen embeddings.  With
\[
 \chi_\lambda(a)=\prod_{\tau,i}\tau(a_i)^{h_{i,\tau}-i},
 \qquad R_a=\chi_\lambda(a)^{-1}\xi_\lambda(a),            \tag{AP6}
\]
\(R_a\) is integral on the original positive cone, and
\(R_z\), for \(z=\operatorname{diag}(\varpi^{n-1},\ldots,1)\),
contracts all nonlowest coordinates into \(\mathfrak m_E\).
The lowest row is fixed by the upper unipotent and has torus
character \(\chi_\lambda\).  Deep compact congruence subgroups act
as one to arbitrary integral precision, first in \textup{(AP5)} and
then in its saturated intersection.  This constructs all the original
weight-lattice requirements without division by a symmetrizer.  Retain
the actual coefficient rank and the original tame type lattice.

\emph{An ordinary vector.}
Put \(\phi=x\theta/\chi_\lambda\).  The original torus equation gives
\(\phi_i(u)=\omega_i(u)^{-1}\) on units.  At a finite depth killing
these characters, evaluation at one in \textup{(AP4)} is a nonzero joint
eigenfunctional on the finite-dimensional congruence-fixed space:
a function supported on \(B(k)I_1(c)\) has value one there.
The torus and positive double-coset operators commute.  Their simultaneous
generalized eigenspace detected by evaluation is nonzero; the commuting
nilpotent parts have a common nonzero kernel.  This gives an actual
joint eigenvector, without assuming semisimplicity of ramified Hecke
operators.

For \(a_j=\operatorname{diag}(\varpi,\ldots,\varpi,1,\ldots,1)\),
with \(j\) copies of \(\varpi\), the unnormalized double coset has
\(q^{j(n-j)}\) terms.  Its Borel character and half modulus give
\[
 q^{j(n-j)}\prod_{i<j}\omega_i(\varpi)q^{-(n-1-i)}
                       =\prod_{i<j}\omega_i(\varpi)q^{-i}.
                                                        \tag{AP7}
\]
After the actual normalization by \(\chi_\lambda(a_j)\), the eigenvalue is
\[
 u_{v,j}=\prod_{i<j}\psi_i(g_\varpi)
              \left(\frac{N_{k/\mathbb Q_\ell}(\varpi)}q\right)^i
          =\prod_{i<j}\psi_i(g_\varpi)\epsilon(g_\varpi)^i.
                                                        \tag{AP8}
\]
It is a unit and gives exactly the original uniformizer equation.
The factor \(N(\varpi)/q\) is retained: it need not equal one.
Tensor the local vectors with the original tame, spherical and infinity
vectors.  This gives a nonzero global classical eigenvector \(F\).
Its finite-dimensional algebraic eigenspace is defined over a finite
coefficient extension.  Scale its finitely many untwisted coset values
into the lattice \textup{(AP5)}; integral compact actions then make all
untwisted values integral.

\emph{Nonzero evaluation in the original module.}
Define
\[
       \operatorname{ev}(F)(g)=e_0^*\xi_\lambda(g_\ell)^{-1}F(g).
                                                        \tag{AP9}
\]
Its torus character is \(x\theta\); its tame type and levels are the
original ones.  At each finite coefficient precision, deeper opposite
unipotent levels follow from the lattice congruence property.
The lowest-row identities intertwine normalized classical operators
with the original scalar operators, term by term.

If \(\operatorname{ev}(F)=0\), every untwisted value has zero lowest
coordinate.  In each normalized \(U_z\) summand, \(R_z\) contracts
that value and the upper-unipotent matrix is integral.  Thus \(U_z\)
sends the kernel lattice into \(\mathfrak m_E\) times itself.  Its
eigenvalue on \(F\) is a unit by \textup{(AP8)}; iteration and separatedness
force \(F=0\), a contradiction.  Hence evaluation is nonzero.
After scaling, its unit eigenvalue puts it in the original integral
ordinary module over \(\mathcal O_E\).  The established exact coefficient
change identifies this with \(M(A,\theta)\otimes_A\mathcal O_E\).
Frobenius polynomial reduction and \textup{(AP8)} give all original
spherical and ordinary residual eigenvalues, so the prescribed residual
Hecke idempotent acts as one.  This proves \(M_A\ne0\).

\emph{Identification of the point.}
The resulting eigenline gives \(y:T\to E\) on the actual integral
Hecke algebra.  Finiteness over \(A\) and the residual congruences make
\(y\) integral, local and continuous.  At the original dense Frobenius
set, \(yf\) and \(x\) have equal determinant polynomials, hence equal
determinants.  The torus equation and \textup{(AP8)} identify their
ordered characters on units and the original uniformizer, hence on
the whole local Galois group.

These polynomial-law coefficients and character values topologically
generate \(B\): in each finite quotient, the subalgebra they generate
carries the same universal pair, so uniqueness in the universal
property makes it the whole quotient.  The continuous maps \(yf\)
and \(x\) consequently agree.  Therefore \(J=\ker f\subseteq\ker x\).
\end{proof}

For the actual seeds, the exact ordered weights and coefficient weights
are respectively
\[
\begin{array}{c|c|c}
 & h & \lambda\\ \hline
 n=3 &(1-m,1,m+1)&(m-1,0,1-m)\\
 n=2,\ a_\infty=-1 &(1-m,1+m)&(m,1-m)\\
 n=2,\ a_\infty=0 &(-m,m)&(m-1,-m).
\end{array}
\]
Here \(m\geq2\) supplies the noncritical gaps, and the original
prime-to-\(\ell\) inertia argument supplies unramifiedness at \(v_0\).
Thus the arithmetic-point hypothesis in the seed specialization is
now supplied.  Characteristic-prime marked support, two-sided global
rigidity and the uniform integral torsion budget remain separate:
a characteristic-zero eigenvector alone does not prove nonemptiness
of every required characteristic-prime fibre.
